% --- Archon physics formalization source begin ---
% archon:physics
% archon:covers PhyXMiniProblems/problem_phyx_mini_0835.lean
% archon:source-report reports/phyx_mini/problem_phyx_mini_0835.source.json

\chapter{Physics problem phyx\_mini\_0835}
\label{ch:PhyXMiniProblems_problem_phyx_mini_0835}

\paragraph{Problem source.}
\#\# Physical scenario

Figure shows a thin rod of length L with total charge Q.

\#\# Question

Evaluate \textbackslash{}( E \textbackslash{}) at \textbackslash{}( r = 3.0\textbackslash{},\textbackslash{}mathrm\{cm\} \textbackslash{}) if \textbackslash{}( L = 5.0\textbackslash{},\textbackslash{}mathrm\{cm\} \textbackslash{}) and \textbackslash{}( Q = 3.0\textbackslash{},\textbackslash{}mathrm\{nC\} \textbackslash{}).

\#\# Answer choices

- A: A: 1.3 \textbackslash{}times 10\textasciicircum{}\{5\}N/C

- B: B: 5.3 \textbackslash{}times 10\textasciicircum{}\{4\}N/C

- C: C: 1.0 \textbackslash{}times 10\textasciicircum{}\{5\}N/C

- D: D: 9.8 \textbackslash{}times 10\textasciicircum{}\{4\}N/C

\#\# Figure caption (auxiliary; use the image as primary evidence)

The image depicts a rectangular bar with a uniform distribution of positive charges, represented by red plus signs. The bar has a length labeled "L" with arrows indicating this length along the bar.

To the right of the bar is a point labeled "P," indicated by a dot. The distance between the closest end of the bar and point P is labeled "r," with an arrow showing this distance.

The relationships between the elements suggest a standard setup for discussing electric fields, with the charged bar influencing the point P.

\#\# Dataset metadata

Electrostatics; Physical Model Grounding Reasoning, Numerical Reasoning

\paragraph{Recorded answer/context.}
D: D: 9.8 \textbackslash{}times 10\textasciicircum{}\{4\}N/C

\paragraph{Figure/image path.}
/root/proposal\_for\_physic/science-mango/phyx\_mini\_run/phyx\_data/test\_image/835.png

\paragraph{Formalization target.}
create a compiling Lean file with sorry bodies at `PhyXMiniProblems/problem\_phyx\_mini\_0835.lean`. The Lean declarations must preserve the physical quantities, dimensions or dimensional roles, figure labels, governing-law hypotheses, and final relation expressed by this problem.
Use Mathlib/Physlib names found through LeanExplore where available. If a physics API is missing, introduce faithful local abstractions rather than scalar placeholder aliases.

\begin{theorem}[Physics formalization target]
\label{thm:physics:phyx_mini_0835:target}
\lean{PhyXMiniProblems.ProblemPhyXMini0835.chargedRodElectricFieldAtP}
\uses{def:physics:phyx-mini-0835:phyxminiproblems-problemphyxmini0835-electricfieldstrengthinnewtonspercoulomb, def:physics:phyx-mini-0835:phyxminiproblems-problemphyxmini0835-axialchargedrodsetup, def:physics:phyx-mini-0835:phyxminiproblems-problemphyxmini0835-matchesprimarychargedrodfigure, def:physics:phyx-mini-0835:phyxminiproblems-problemphyxmini0835-matcheschargedrodproblemdata, def:physics:phyx-mini-0835:phyxminiproblems-problemphyxmini0835-hasphysicalchargedrodparameters, def:physics:phyx-mini-0835:phyxminiproblems-problemphyxmini0835-usesvacuumcoulombconstant, def:physics:phyx-mini-0835:phyxminiproblems-problemphyxmini0835-satisfiesuniformfinitelinechargemodel, def:physics:phyx-mini-0835:phyxminiproblems-problemphyxmini0835-satisfiesaxialcoulombsuperpositionlaw, def:physics:phyx-mini-0835:phyxminiproblems-problemphyxmini0835-fieldatphasstoredaxialmagnitude, def:physics:phyx-mini-0835:phyxminiproblems-problemphyxmini0835-roundstonearestthousand, def:physics:phyx-mini-0835:phyxminiproblems-problemphyxmini0835-isuniqueroundedanswer}
The assigned autoformalize agent should translate this physics problem into Lean declarations in the covered file, with theorem and lemma proof bodies written as `by sorry`.
\end{theorem}
\begin{proof}
This is an autoformalization task, not a proof task. Produce statements that can later be proved by the physics prover without weakening the physical model.
\end{proof}

% --- Archon declaration topology begin ---
\section{Declaration topology}
\begin{definition}[linear Charge Density Dimension]
  \label{def:physics:phyx-mini-0835:phyxminiproblems-problemphyxmini0835-linearchargedensitydimension}
  \lean{PhyXMiniProblems.ProblemPhyXMini0835.linearChargeDensityDimension}
  The dimension C L⁻¹ of linear electric-charge density
\end{definition}
\begin{proof}
  This is the declared model definition of linear Charge Density Dimension.
\end{proof}
\begin{definition}[electric Field Strength Dimension]
  \label{def:physics:phyx-mini-0835:phyxminiproblems-problemphyxmini0835-electricfieldstrengthdimension}
  \lean{PhyXMiniProblems.ProblemPhyXMini0835.electricFieldStrengthDimension}
  The dimension M L T⁻² C⁻¹ of electric-field strength
\end{definition}
\begin{proof}
  This is the declared model definition of electric Field Strength Dimension.
\end{proof}
\begin{definition}[coulomb Constant Dimension]
  \label{def:physics:phyx-mini-0835:phyxminiproblems-problemphyxmini0835-coulombconstantdimension}
  \lean{PhyXMiniProblems.ProblemPhyXMini0835.coulombConstantDimension}
  The dimension M L³ T⁻² C⁻² of Coulomb's constant
\end{definition}
\begin{proof}
  This is the declared model definition of coulomb Constant Dimension.
\end{proof}
\begin{definition}[Length Quantity]
  \label{def:physics:phyx-mini-0835:phyxminiproblems-problemphyxmini0835-lengthquantity}
  \lean{PhyXMiniProblems.ProblemPhyXMini0835.LengthQuantity}
  A nonnegative, unit-independent physical length
\end{definition}
\begin{proof}
  This is the declared model definition of Length Quantity.
\end{proof}
\begin{definition}[Charge Magnitude Quantity]
  \label{def:physics:phyx-mini-0835:phyxminiproblems-problemphyxmini0835-chargemagnitudequantity}
  \lean{PhyXMiniProblems.ProblemPhyXMini0835.ChargeMagnitudeQuantity}
  A nonnegative, unit-independent electric-charge magnitude
\end{definition}
\begin{proof}
  This is the declared model definition of Charge Magnitude Quantity.
\end{proof}
\begin{definition}[Linear Charge Density Quantity]
  \label{def:physics:phyx-mini-0835:phyxminiproblems-problemphyxmini0835-linearchargedensityquantity}
  \lean{PhyXMiniProblems.ProblemPhyXMini0835.LinearChargeDensityQuantity}
  \uses{def:physics:phyx-mini-0835:phyxminiproblems-problemphyxmini0835-linearchargedensitydimension}
  A nonnegative, unit-independent linear charge density
\end{definition}
\begin{proof}
  This is the declared model definition of Linear Charge Density Quantity.
\end{proof}
\begin{definition}[Electric Field Strength Quantity]
  \label{def:physics:phyx-mini-0835:phyxminiproblems-problemphyxmini0835-electricfieldstrengthquantity}
  \lean{PhyXMiniProblems.ProblemPhyXMini0835.ElectricFieldStrengthQuantity}
  \uses{def:physics:phyx-mini-0835:phyxminiproblems-problemphyxmini0835-electricfieldstrengthdimension}
  A nonnegative, unit-independent electric-field magnitude
\end{definition}
\begin{proof}
  This is the declared model definition of Electric Field Strength Quantity.
\end{proof}
\begin{definition}[Coulomb Constant Quantity]
  \label{def:physics:phyx-mini-0835:phyxminiproblems-problemphyxmini0835-coulombconstantquantity}
  \lean{PhyXMiniProblems.ProblemPhyXMini0835.CoulombConstantQuantity}
  \uses{def:physics:phyx-mini-0835:phyxminiproblems-problemphyxmini0835-coulombconstantdimension}
  A nonnegative, dimensionful value of Coulomb's constant
\end{definition}
\begin{proof}
  This is the declared model definition of Coulomb Constant Quantity.
\end{proof}
\begin{definition}[length Readout]
  \label{def:physics:phyx-mini-0835:phyxminiproblems-problemphyxmini0835-lengthreadout}
  \lean{PhyXMiniProblems.ProblemPhyXMini0835.lengthReadout}
  \uses{def:physics:phyx-mini-0835:phyxminiproblems-problemphyxmini0835-lengthquantity}
  Read a physical length in a selected Physlib length unit
\end{definition}
\begin{proof}
  This is the declared model definition of length Readout.
\end{proof}
\begin{definition}[length In Meters]
  \label{def:physics:phyx-mini-0835:phyxminiproblems-problemphyxmini0835-lengthinmeters}
  \lean{PhyXMiniProblems.ProblemPhyXMini0835.lengthInMeters}
  \uses{def:physics:phyx-mini-0835:phyxminiproblems-problemphyxmini0835-lengthquantity, def:physics:phyx-mini-0835:phyxminiproblems-problemphyxmini0835-lengthreadout}
  SI metre readout of a physical length
\end{definition}
\begin{proof}
  This is the declared model definition of length In Meters.
\end{proof}
\begin{definition}[length In Centimeters]
  \label{def:physics:phyx-mini-0835:phyxminiproblems-problemphyxmini0835-lengthincentimeters}
  \lean{PhyXMiniProblems.ProblemPhyXMini0835.lengthInCentimeters}
  \uses{def:physics:phyx-mini-0835:phyxminiproblems-problemphyxmini0835-lengthquantity, def:physics:phyx-mini-0835:phyxminiproblems-problemphyxmini0835-lengthreadout}
  Centimetre readout used for the two printed length data
\end{definition}
\begin{proof}
  This is the declared model definition of length In Centimeters.
\end{proof}
\begin{definition}[charge In Coulombs]
  \label{def:physics:phyx-mini-0835:phyxminiproblems-problemphyxmini0835-chargeincoulombs}
  \lean{PhyXMiniProblems.ProblemPhyXMini0835.chargeInCoulombs}
  \uses{def:physics:phyx-mini-0835:phyxminiproblems-problemphyxmini0835-chargemagnitudequantity}
  SI coulomb readout of a physical charge magnitude
\end{definition}
\begin{proof}
  This is the declared model definition of charge In Coulombs.
\end{proof}
\begin{definition}[charge In Nanocoulombs]
  \label{def:physics:phyx-mini-0835:phyxminiproblems-problemphyxmini0835-chargeinnanocoulombs}
  \lean{PhyXMiniProblems.ProblemPhyXMini0835.chargeInNanocoulombs}
  \uses{def:physics:phyx-mini-0835:phyxminiproblems-problemphyxmini0835-chargemagnitudequantity, def:physics:phyx-mini-0835:phyxminiproblems-problemphyxmini0835-chargeincoulombs}
  Nanocoulomb readout used by the printed total charge
\end{definition}
\begin{proof}
  This is the declared model definition of charge In Nanocoulombs.
\end{proof}
\begin{definition}[linear Charge Density In Coulombs Per Meter]
  \label{def:physics:phyx-mini-0835:phyxminiproblems-problemphyxmini0835-linearchargedensityincoulombspermeter}
  \lean{PhyXMiniProblems.ProblemPhyXMini0835.linearChargeDensityInCoulombsPerMeter}
  \uses{def:physics:phyx-mini-0835:phyxminiproblems-problemphyxmini0835-linearchargedensityquantity}
  SI readout of linear charge density, in coulombs per metre
\end{definition}
\begin{proof}
  This is the declared model definition of linear Charge Density In Coulombs Per Meter.
\end{proof}
\begin{definition}[electric Field Strength In Newtons Per Coulomb]
  \label{def:physics:phyx-mini-0835:phyxminiproblems-problemphyxmini0835-electricfieldstrengthinnewtonspercoulomb}
  \lean{PhyXMiniProblems.ProblemPhyXMini0835.electricFieldStrengthInNewtonsPerCoulomb}
  \uses{def:physics:phyx-mini-0835:phyxminiproblems-problemphyxmini0835-electricfieldstrengthquantity}
  SI readout of electric-field strength, in newtons per coulomb
\end{definition}
\begin{proof}
  This is the declared model definition of electric Field Strength In Newtons Per Coulomb.
\end{proof}
\begin{definition}[coulomb Constant In Newton Square Meters Per Coulomb Squared]
  \label{def:physics:phyx-mini-0835:phyxminiproblems-problemphyxmini0835-coulombconstantinnewtonsquaremeterspercoulombsquared}
  \lean{PhyXMiniProblems.ProblemPhyXMini0835.coulombConstantInNewtonSquareMetersPerCoulombSquared}
  \uses{def:physics:phyx-mini-0835:phyxminiproblems-problemphyxmini0835-coulombconstantquantity}
  SI readout of Coulomb's constant, in newton-square-metres per coulomb squared
\end{definition}
\begin{proof}
  This is the declared model definition of coulomb Constant In Newton Square Meters Per Coulomb Squared.
\end{proof}
\begin{definition}[Thin Uniform Charged Rod]
  \label{def:physics:phyx-mini-0835:phyxminiproblems-problemphyxmini0835-thinuniformchargedrod}
  \lean{PhyXMiniProblems.ProblemPhyXMini0835.ThinUniformChargedRod}
  \uses{def:physics:phyx-mini-0835:phyxminiproblems-problemphyxmini0835-lengthquantity, def:physics:phyx-mini-0835:phyxminiproblems-problemphyxmini0835-chargemagnitudequantity, def:physics:phyx-mini-0835:phyxminiproblems-problemphyxmini0835-linearchargedensityquantity}
  The thin rod together with its independent charge observables
\end{definition}
\begin{proof}
  This is the declared model definition of Thin Uniform Charged Rod.
\end{proof}
\begin{definition}[Rod Axis Marker]
  \label{def:physics:phyx-mini-0835:phyxminiproblems-problemphyxmini0835-rodaxismarker}
  \lean{PhyXMiniProblems.ProblemPhyXMini0835.RodAxisMarker}
  Named axial markers appearing in, or geometrically determined by, the image
\end{definition}
\begin{proof}
  This is the declared model definition of Rod Axis Marker.
\end{proof}
\begin{definition}[Figure Feature]
  \label{def:physics:phyx-mini-0835:phyxminiproblems-problemphyxmini0835-figurefeature}
  \lean{PhyXMiniProblems.ProblemPhyXMini0835.FigureFeature}
  Features explicitly visible in the supplied image 835.png
\end{definition}
\begin{proof}
  This is the declared model definition of Figure Feature.
\end{proof}
\begin{definition}[Charged Rod Figure]
  \label{def:physics:phyx-mini-0835:phyxminiproblems-problemphyxmini0835-chargedrodfigure}
  \lean{PhyXMiniProblems.ProblemPhyXMini0835.ChargedRodFigure}
  \uses{def:physics:phyx-mini-0835:phyxminiproblems-problemphyxmini0835-lengthquantity, def:physics:phyx-mini-0835:phyxminiproblems-problemphyxmini0835-rodaxismarker, def:physics:phyx-mini-0835:phyxminiproblems-problemphyxmini0835-figurefeature}
  The axial coordinates are explicit metre readouts of a schematic coordinate chart. The two label quantities themselves remain dimensionful
\end{definition}
\begin{proof}
  This is the declared model definition of Charged Rod Figure.
\end{proof}
\begin{definition}[Axial Charged Rod Setup]
  \label{def:physics:phyx-mini-0835:phyxminiproblems-problemphyxmini0835-axialchargedrodsetup}
  \lean{PhyXMiniProblems.ProblemPhyXMini0835.AxialChargedRodSetup}
  \uses{def:physics:phyx-mini-0835:phyxminiproblems-problemphyxmini0835-lengthquantity, def:physics:phyx-mini-0835:phyxminiproblems-problemphyxmini0835-electricfieldstrengthquantity, def:physics:phyx-mini-0835:phyxminiproblems-problemphyxmini0835-coulombconstantquantity, def:physics:phyx-mini-0835:phyxminiproblems-problemphyxmini0835-thinuniformchargedrod, def:physics:phyx-mini-0835:phyxminiproblems-problemphyxmini0835-chargedrodfigure}
  The independent physical apparatus and observables. In particular, neither the electric-field magnitude nor the vector field is defined from a displayed answer choice
\end{definition}
\begin{proof}
  This is the declared model definition of Axial Charged Rod Setup.
\end{proof}
\begin{definition}[Matches Primary Charged Rod Figure]
  \label{def:physics:phyx-mini-0835:phyxminiproblems-problemphyxmini0835-matchesprimarychargedrodfigure}
  \lean{PhyXMiniProblems.ProblemPhyXMini0835.MatchesPrimaryChargedRodFigure}
  \uses{def:physics:phyx-mini-0835:phyxminiproblems-problemphyxmini0835-lengthinmeters, def:physics:phyx-mini-0835:phyxminiproblems-problemphyxmini0835-axialchargedrodsetup}
  Primary-image transcription. Crucially, the r arrow starts at the rod center and ends at P; no electric-field value or answer label occurs here
\end{definition}
\begin{proof}
  This is the declared model definition of Matches Primary Charged Rod Figure.
\end{proof}
\begin{definition}[Matches Charged Rod Problem Data]
  \label{def:physics:phyx-mini-0835:phyxminiproblems-problemphyxmini0835-matcheschargedrodproblemdata}
  \lean{PhyXMiniProblems.ProblemPhyXMini0835.MatchesChargedRodProblemData}
  \uses{def:physics:phyx-mini-0835:phyxminiproblems-problemphyxmini0835-lengthinmeters, def:physics:phyx-mini-0835:phyxminiproblems-problemphyxmini0835-lengthincentimeters, def:physics:phyx-mini-0835:phyxminiproblems-problemphyxmini0835-chargeincoulombs, def:physics:phyx-mini-0835:phyxminiproblems-problemphyxmini0835-chargeinnanocoulombs, def:physics:phyx-mini-0835:phyxminiproblems-problemphyxmini0835-axialchargedrodsetup}
  Numerical data stated in the question, retaining both printed and SI readouts
\end{definition}
\begin{proof}
  This is the declared model definition of Matches Charged Rod Problem Data.
\end{proof}
\begin{definition}[Has Physical Charged Rod Parameters]
  \label{def:physics:phyx-mini-0835:phyxminiproblems-problemphyxmini0835-hasphysicalchargedrodparameters}
  \lean{PhyXMiniProblems.ProblemPhyXMini0835.HasPhysicalChargedRodParameters}
  \uses{def:physics:phyx-mini-0835:phyxminiproblems-problemphyxmini0835-lengthinmeters, def:physics:phyx-mini-0835:phyxminiproblems-problemphyxmini0835-chargeincoulombs, def:physics:phyx-mini-0835:phyxminiproblems-problemphyxmini0835-linearchargedensityincoulombspermeter, def:physics:phyx-mini-0835:phyxminiproblems-problemphyxmini0835-coulombconstantinnewtonsquaremeterspercoulombsquared, def:physics:phyx-mini-0835:phyxminiproblems-problemphyxmini0835-axialchargedrodsetup}
  Positivity and separation assumptions selecting the depicted external point
\end{definition}
\begin{proof}
  This is the declared model definition of Has Physical Charged Rod Parameters.
\end{proof}
\begin{definition}[Uses Vacuum Coulomb Constant]
  \label{def:physics:phyx-mini-0835:phyxminiproblems-problemphyxmini0835-usesvacuumcoulombconstant}
  \lean{PhyXMiniProblems.ProblemPhyXMini0835.UsesVacuumCoulombConstant}
  \uses{def:physics:phyx-mini-0835:phyxminiproblems-problemphyxmini0835-coulombconstantinnewtonsquaremeterspercoulombsquared, def:physics:phyx-mini-0835:phyxminiproblems-problemphyxmini0835-axialchargedrodsetup}
  The dimensionful Coulomb constant agrees with Physlib's electromagnetic system and has the standard vacuum SI calibration. This is reference data, not the requested field value
\end{definition}
\begin{proof}
  This is the declared model definition of Uses Vacuum Coulomb Constant.
\end{proof}
\begin{definition}[Satisfies Uniform Finite Line Charge Model]
  \label{def:physics:phyx-mini-0835:phyxminiproblems-problemphyxmini0835-satisfiesuniformfinitelinechargemodel}
  \lean{PhyXMiniProblems.ProblemPhyXMini0835.SatisfiesUniformFiniteLineChargeModel}
  \uses{def:physics:phyx-mini-0835:phyxminiproblems-problemphyxmini0835-lengthinmeters, def:physics:phyx-mini-0835:phyxminiproblems-problemphyxmini0835-chargeincoulombs, def:physics:phyx-mini-0835:phyxminiproblems-problemphyxmini0835-linearchargedensityincoulombspermeter, def:physics:phyx-mini-0835:phyxminiproblems-problemphyxmini0835-axialchargedrodsetup}
  A uniform finite line charge has constant density Q/L, and integrating that density over its centered support recovers Q. Neither law mentions the electric field or any answer choice
\end{definition}
\begin{proof}
  This is the declared model definition of Satisfies Uniform Finite Line Charge Model.
\end{proof}
\begin{definition}[Satisfies Axial Coulomb Superposition Law]
  \label{def:physics:phyx-mini-0835:phyxminiproblems-problemphyxmini0835-satisfiesaxialcoulombsuperpositionlaw}
  \lean{PhyXMiniProblems.ProblemPhyXMini0835.SatisfiesAxialCoulombSuperpositionLaw}
  \uses{def:physics:phyx-mini-0835:phyxminiproblems-problemphyxmini0835-lengthinmeters, def:physics:phyx-mini-0835:phyxminiproblems-problemphyxmini0835-linearchargedensityincoulombspermeter, def:physics:phyx-mini-0835:phyxminiproblems-problemphyxmini0835-electricfieldstrengthinnewtonspercoulomb, def:physics:phyx-mini-0835:phyxminiproblems-problemphyxmini0835-coulombconstantinnewtonsquaremeterspercoulombsquared, def:physics:phyx-mini-0835:phyxminiproblems-problemphyxmini0835-axialchargedrodsetup}
  Coulomb superposition for a point P on the positive rod axis. A source element at axial coordinate x is r - x metres from P, so its field contribution has magnitude k λ dx / (r - x)². This integral law is general for the stated rod data and does not assume the requested closed form or numerical result
\end{definition}
\begin{proof}
  This is the declared model definition of Satisfies Axial Coulomb Superposition Law.
\end{proof}
\begin{definition}[Field At PHas Stored Axial Magnitude]
  \label{def:physics:phyx-mini-0835:phyxminiproblems-problemphyxmini0835-fieldatphasstoredaxialmagnitude}
  \lean{PhyXMiniProblems.ProblemPhyXMini0835.FieldAtPHasStoredAxialMagnitude}
  \uses{def:physics:phyx-mini-0835:phyxminiproblems-problemphyxmini0835-electricfieldstrengthinnewtonspercoulomb, def:physics:phyx-mini-0835:phyxminiproblems-problemphyxmini0835-axialchargedrodsetup}
  The positive rod produces an electric field in the positive axial direction at P. This relates the dimensionful magnitude to Physlib's vector-field object without fixing the magnitude numerically
\end{definition}
\begin{proof}
  This is the declared model definition of Field At PHas Stored Axial Magnitude.
\end{proof}
\begin{lemma}[axial Electric Field closed Form]
  \label{lem:physics:phyx-mini-0835:phyxminiproblems-problemphyxmini0835-axialelectricfield-closedform}
  \lean{PhyXMiniProblems.ProblemPhyXMini0835.axialElectricField_closedForm}
  \uses{def:physics:phyx-mini-0835:phyxminiproblems-problemphyxmini0835-lengthinmeters, def:physics:phyx-mini-0835:phyxminiproblems-problemphyxmini0835-chargeincoulombs, def:physics:phyx-mini-0835:phyxminiproblems-problemphyxmini0835-electricfieldstrengthinnewtonspercoulomb, def:physics:phyx-mini-0835:phyxminiproblems-problemphyxmini0835-coulombconstantinnewtonsquaremeterspercoulombsquared, def:physics:phyx-mini-0835:phyxminiproblems-problemphyxmini0835-axialchargedrodsetup, def:physics:phyx-mini-0835:phyxminiproblems-problemphyxmini0835-hasphysicalchargedrodparameters, def:physics:phyx-mini-0835:phyxminiproblems-problemphyxmini0835-satisfiesuniformfinitelinechargemodel, def:physics:phyx-mini-0835:phyxminiproblems-problemphyxmini0835-satisfiesaxialcoulombsuperpositionlaw}
  Integrating the axial Coulomb law for the uniform rod yields E = k Q / (r² - (L/2)²). This is a derived lemma, not a premise of the main numerical theorem
\end{lemma}
\begin{proof}
  The conclusion follows from the governing relations and figure readouts named in the statement of axial Electric Field closed Form.
\end{proof}
\begin{definition}[Answer Choice]
  \label{def:physics:phyx-mini-0835:phyxminiproblems-problemphyxmini0835-answerchoice}
  \lean{PhyXMiniProblems.ProblemPhyXMini0835.AnswerChoice}
  Labels of the four field-strength answers printed in the question
\end{definition}
\begin{proof}
  This is the declared model definition of Answer Choice.
\end{proof}
\begin{definition}[answer Field Strength In Newtons Per Coulomb]
  \label{def:physics:phyx-mini-0835:phyxminiproblems-problemphyxmini0835-answerfieldstrengthinnewtonspercoulomb}
  \lean{PhyXMiniProblems.ProblemPhyXMini0835.answerFieldStrengthInNewtonsPerCoulomb}
  \uses{def:physics:phyx-mini-0835:phyxminiproblems-problemphyxmini0835-answerchoice}
  Printed field magnitude, in newtons per coulomb, beside each choice
\end{definition}
\begin{proof}
  This is the declared model definition of answer Field Strength In Newtons Per Coulomb.
\end{proof}
\begin{definition}[Rounds To Nearest Thousand]
  \label{def:physics:phyx-mini-0835:phyxminiproblems-problemphyxmini0835-roundstonearestthousand}
  \lean{PhyXMiniProblems.ProblemPhyXMini0835.RoundsToNearestThousand}
  The answers are printed to the nearest 10\textasciicircum{}3 N/C. Strictly lying within half that increment faithfully expresses rounding to a displayed value and does not define the field itself
\end{definition}
\begin{proof}
  This is the declared model definition of Rounds To Nearest Thousand.
\end{proof}
\begin{definition}[Is Unique Rounded Answer]
  \label{def:physics:phyx-mini-0835:phyxminiproblems-problemphyxmini0835-isuniqueroundedanswer}
  \lean{PhyXMiniProblems.ProblemPhyXMini0835.IsUniqueRoundedAnswer}
  \uses{def:physics:phyx-mini-0835:phyxminiproblems-problemphyxmini0835-answerchoice, def:physics:phyx-mini-0835:phyxminiproblems-problemphyxmini0835-answerfieldstrengthinnewtonspercoulomb, def:physics:phyx-mini-0835:phyxminiproblems-problemphyxmini0835-roundstonearestthousand}
  A displayed choice is the unique one to which the computed field rounds
\end{definition}
\begin{proof}
  This is the declared model definition of Is Unique Rounded Answer.
\end{proof}
% --- Archon declaration topology end ---
% --- Archon physics formalization source end ---
